\section{Баєсова регресія}

У баєсовій моделі коефіцієнти трактуються як %
випадкові величини з апріорним розподілом. Після спостереження навчальних даних %
отримується апостеріорний розподіл, %
математичне сподівання якого використовується для %
прогнозування кольору цільового регіону.

У даному алгоритмі машинного навчання лежить в основі вищеописана лінійна регресія.
Тут так само визначена регресійна модель \eqref{eq:formula-5.4}:
\begin{equation}
	y_{i} = \alpha + \beta x_{i} + \varepsilon_{i}, \quad i = 1, \dots, n. \label{eq:formula-5.4}
\end{equation}

Тут $\varepsilon_{i}$ --- незалежні та нормально розподілені випадкові величини з
нульовим середнім та дисперсією $\sigma^{2}$.

Визначимо функцію максимальної правдоподібності \eqref{eq:formula-5.5}:
\begin{equation}
	\rho(y \mid X, \beta, \sigma^{2}) = (\sigma^{2})^{-\frac{n}{2}}\exp\left(-\frac{1}{2\sigma^{2}}
	(y - X\beta)^{T}(y - X\beta)\right) \label{eq:formula-5.5}
\end{equation}

Для того, щоб знайти розв'язок \eqref{eq:formula-5.6} за допомогою відомого
методу найменших квадратів використовують псевдообернення Мура-Пенроуза:
\begin{equation}
	\hat{\beta}= (X^{T} X)^{-1}X^{T} y \label{eq:formula-5.6}
\end{equation}

де $X$ --- матриця із розмірністю $n \times k$, причому кожен із рядків представлений
у вигляді $x_{i}^{T}$, а $y$ є вектором $[y_{1}, y_{2}, \dots, y_{n}]^{T}$.